\subsection{无处不在噪声}
噪声存在于各个前后端的各个环节，按产生的逻辑可以分为系统噪声和观测噪声。

\noindent\textbf{系统噪声：}
\begin{itemize}
    \item 运动噪声：机器人在赛场上运动时候的抖动、震动，属于被观测的系统
    \item 算法噪声：由于各种原因，如离散化等处理手段导致的数据异常（跳变经常会看到），属于观测或预测的系统
\end{itemize}
\textbf{观测噪声}
\begin{itemize}
    \item 传感器噪声：相机传感器的噪声，量化误差等等
    \item 环境噪声：光照变化、反射镜面等之余视觉图像，IMU漂移之余位置估计
\end{itemize}

\subsection{滤波的基本概念}
通过传感器得到的测量值，会包括信号的真实值和无处不在噪声。滤波就是分离有用的信号和噪声，提高系统的信噪比。
根据滤波方式，常被分为时域滤波和频域滤波，这里就先不做展开。

滤波器常见的应用有：
\begin{itemize}
    \item 平滑信号，降低高频噪声：如取出抖动毛刺等，也如电容的作用
    \item 增强特定频率的成分：如保留目标特征，像放大器这样的
    \item 进行边缘检测、特征提取：像opencv里面学的各种图像算法
    \item 优化传感器数据：如IMU的数据融合，消除抖动的噪声
\end{itemize}

\subsection{常见滤波器及其原理}
\subsubsection{低通滤波器}
对时域信号进行操作，平滑信号，去除高频噪声。如图像里的随机噪声，IMU的高速震动等。
$$Y(n)=aX(n)+(1-a)Y(n-1)$$
其中a为滤波系数，控制平滑程度，a=1时无效果

\textbf{实现方法：}
\begin{itemize}
    \item 移动平均：移动窗口算法
    \item 高斯滤波：使用高斯权重
    \item 指数平滑滤波
\end{itemize}

\textbf{代码实现}


\subsubsection{高通滤波器}
取出低频分量，保留高频信息，常应用在边缘检测。

\subsection{卡尔曼滤波}
卡尔曼滤波（Kalman filter）是一种高效的时域递归滤波器，能够从不完全且含有噪声的测量数据中，对动态系统的状态进行最优估计。

卡尔曼滤波的核心思想是利用\textbf{预测（Predict）}和\textbf{更新（Update）}两个步骤来估计系统的状态：
\begin{itemize}
    \item \textbf{预测阶段：}基于上一时刻的状态，利用系统的数学模型预测当前时刻的状态，并给出不确定性。
    \item \textbf{更新阶段：}在新的观测数据到来后，使用加权平均的方法修正预测值，其中\textbf{置信度越高的数据}权重越大，从而得到更精确的估计。
\end{itemize}
简单来说，卡尔曼滤波就是对\textbf{测量值}和\textbf{预测值}进行加权平均，以减少噪声的影响。

卡尔曼滤波在许多领域具有广泛应用，例如：
\begin{itemize}
    \item 目标追踪：在识别装甲板的自瞄算法和击打能量机关有重要应用
    \item 机器人定位：融合IMU，里程计等信息（信息融合）
    \item 姿态估计：进行陀螺仪数据解算
\end{itemize}

\subsubsection{公式推导}
抛开令人头大的定义，我们可以从数据融合的角度来理解卡尔曼滤波的推导过程。它的本质是融合先验状态估计和测量数据，并通过计算\textbf{卡尔曼增益（Kalman Gain）}来确定融合的最佳权重。

\begin{enumerate}
    \item \textbf{卡尔曼滤波的预测}
    
    过程基于动态系统模型，它可以描述如下状态空间方程：
    \begin{equation*}
        \mathbf{x}_k = \mathbf{A} \mathbf{x}_{k-1} + \mathbf{B}_k \mathbf{u}_k + \mathbf{w}_k
    \end{equation*}
    
    其中：
    
    $\mathbf{x}_k$：系统在 $k$ 时刻的状态变量，例如位置、速度、加速度等。
    
    $\mathbf{A}$：状态转移矩阵，描述状态变量在时间上的演变。
    
    $\mathbf{B}_k$：控制矩阵（若有控制输入）。
    
    $\mathbf{u}_k$：控制输入（例如外力影响）。
    
    $\mathbf{w}_k$：过程噪声（系统内部的不确定性）。

    该公式表示：当前状态由上一时刻状态加上某些外部影响项组成。

    \item \textbf{测量模型（观测）}

    系统的测量数据通常与实际状态之间存在某种线性关系，并伴随测量噪声，描述如下：
    \begin{equation}
        \mathbf{z}_k = \mathbf{H}_k \mathbf{x}_k + \mathbf{v}_k
    \end{equation}

    其中：

    $\mathbf{z}_k$：观测值，即传感器测量得到的变量。

    $\mathbf{H}_k$：观测矩阵，描述状态变量如何映射到观测值。

    $\mathbf{v}_k$：测量噪声（例如传感器误差）。

    这表示：我们不能直接测量真实状态，而是通过传感器间接测量并受到噪声干扰。

    \item \textbf{数据融合与卡尔曼增益}

    两个方程都涉及状态变量 $\mathbf{x}_k$，我们需要通过数据融合来求得最优估计值。目标是找到卡尔曼增益，使得融合后的估计值与实际状态的误差方差最小。

    卡尔曼增益的计算如下：
    \begin{equation}
        \mathbf{K}_k = \frac{\mathbf{P}_{k}^- \mathbf{H}^T}{\mathbf{H} \mathbf{P}_{k}^- \mathbf{H}^T + \mathbf{R}}
    \end{equation}

    最终的最优状态估计值：
    \begin{equation}
        \mathbf{\hat{x}_{k}} = \mathbf{\hat{x}_{k}} + \mathbf{K}_k (\mathbf{Z}_k - \mathbf{H} \mathbf{\hat{x}_k^-})
    \end{equation}

    其中：

    $\mathbf{K}_k$：卡尔曼增益，决定了预测值和测量值的加权比例。

    $\mathbf{P}_{k}^-$：上一时刻的状态协方差矩阵，表示预测的不确定性。

    $\mathbf{R}_k$：测量噪声协方差，表示测量值的可信度。

\end{enumerate}

\textbf{直观理解}
\begin{itemize}
    \item 如果 $\mathbf{R}_k$ 很小（测量值可靠），则卡尔曼增益 $\mathbf{K}_k$ 增大，意味着测量值权重更高，最终状态更接近测量值。
    \item 如果 $\mathbf{R}_k$ 很大（测量值不可靠），则 $\mathbf{K}_k$ 减小，意味着更依赖预测值，最终状态更接近先验估计。
\end{itemize}

卡尔曼滤波本质上是一种动态数据融合方法，它在预测和测量之间找到一个平衡点，使得最终的状态估计值更接近真实值。它广泛用于目标跟踪、姿态估计、机器人定位等应用中，极大提高了测量精度。

\subsubsection{算法实现}
实现分为两步，一个是预测，一个是更新。
预测部分就是获得先验的状态，和先验估计和真实值的误差的协方差。
更新即根据先验状态方差更新卡尔曼增益，使用卡尔曼增益来融合数据，得到估计的状态，再去更新估计值和真实值的误差的协方差。

\subsubsection{拓展部分}
拓展卡尔曼
轨迹卡尔曼

\subsubsection{卡尔曼滤波实战}